\begin{enumerate}
  \item 干扰源在一次测试中位置、频道、有效接收半径和定向方向保持不变，各频道至多对应一个未清除源。该条件来自题面和模拟器协议，是按频道独立维护定位区域的依据。
  \item 示向误差只采用题面给出的 $\pm1^\circ$ 硬边界，不假设其独立、均匀或均值为零；同点重复检测不增加信息。
  \item 首次返回 \texttt{direction} 时，源点到检测点距离大于 $5\unit{m}$ 且不超过 $1500\unit{m}$。为获得保证接收结论，第二检测点设计统一使用 $1000\unit{m}$ 的接收半径下界。
  \item 机器狗可按题面在平面内直线移动，不另设道路、障碍或圆域行动边界。圆域只约束源的位置，覆盖点允许落在圆域外侧。
  \item 数值计算使用双精度。示向楔形半角在机器狗实现中取 $1.01^\circ$，其中 $0.01^\circ$ 用于覆盖接口两位小数舍入与浮点误差；论文解析推导使用题面值 $1^\circ$。
  \item 合成模拟中源位置、接收半径与定向概率的分布仅用于压力测试，不作为官方案例的概率模型，也不据此推断正式成绩。
\end{enumerate}
